\section{Prime-layer radical bounds and the corrected Mersenne order ledger}
\label{sec:mersenne-prime-layer}

Put $M_m=2^m-1$ and
\[
                 W_m=\frac{M_m}{\rad(M_m)}.
\]
For every odd prime $q$, let $d_q=\operatorname{ord}_q(2)$ and
$w_q=v_q(2^{d_q}-1)$.  Define the base excess of the exact-order block $d$ by
\[
                E_d=\prod_{\substack{q\ {\mathrm{prime}}\\d_q=d}}
                         q^{w_q-1}.
\]
The index-lifting contribution must be kept separate:
\[
 B_m=\prod_{d\mid m}E_d,\qquad
 L_m=\prod_{\substack{q\ {\mathrm{prime}}\\q\mid M_m}}
                         q^{v_q(m/d_q)}.
\]

\begin{theorem}[Correct exact-order decomposition]
\label{thm:mersenne-order-decomposition}
For every $m\ge1$,
\begin{equation}\label{eq:mersenne-order-decomposition}
                       W_m=L_mB_m=L_m\prod_{d\mid m}E_d.
\end{equation}
Moreover $L_m\mid m$, so $\log L_m\le\log m=o(m)$.
\end{theorem}

\begin{proof}
If $q\mid M_m$, then $d_q\mid m$, and odd-prime LTE gives
\[
 v_q(M_m)=w_q+v_q(m/d_q).
\]
After one copy of $q$ is removed into the radical, the first summand supplies
$q^{w_q-1}$ in the unique block $E_{d_q}$ and the second supplies its factor
in $L_m$.  Multiplying over the prime support proves
\eqref{eq:mersenne-order-decomposition}.  Since $d_q\mid q-1$, one has
$q\nmid d_q$ and therefore $v_q(m/d_q)=v_q(m)$.  Thus $L_m$ takes from $m$
only the full powers of a subset of its prime divisors, proving $L_m\mid m$.
\end{proof}

The factor $L_m$ cannot be deleted.  At $m=6$,
\[
 M_6=63=3^2\cdot7,\qquad \rad(M_6)=21,\qquad W_6=3.
\]
Here $d_3=2$, $w_3=1$, $d_7=3$, and $w_7=1$, so the base product is one,
whereas $L_6=3$.  This is a complete counterexample to the formula
$W_m=\prod_{d\mid m}E_d$ without an index-lifting term.  It supports the
corrected decomposition rather than refuting the order-block route.

At a prime index the lifting factor disappears, and the radical admits an
elementary uniform lower bound.

\begin{theorem}[Composite prime-layer square carrier]
\label{thm:mersenne-prime-square-carrier}
Let $\ell$ be an odd prime.  If $M_\ell$ is composite, then it has at least
two distinct prime divisors, every prime divisor $q$ satisfies
\[
 \operatorname{ord}_q(2)=\ell,\qquad
 2\ell\mid q-1,\qquad q\ge2\ell+1,
\]
and consequently
\begin{equation}\label{eq:mersenne-square-carrier}
 (2\ell+1)^2\le\rad(M_\ell),\qquad
 (2\ell+1)^2E_\ell\le\Phi_\ell(2).
\end{equation}
If one prime factor is at least $H$, the square in
\eqref{eq:mersenne-square-carrier} may be replaced by
$H\,(2\ell+1)$.
\end{theorem}

\begin{proof}
First, $2^\ell-1$ is not a nontrivial prime power.  Indeed, if
$2^\ell-1=q^v$, then an even $v$ contradicts reduction modulo eight.  If
$v\ge3$ is odd, the factorization
\[
 q^v+1=(q+1)(q^{v-1}-q^{v-2}+\cdots-q+1)
\]
has an odd second factor strictly greater than one, impossible for a power of
two.  Hence a composite value has two distinct prime factors.

For a prime divisor $q$, its order divides the prime $\ell$ and is not one,
so it equals $\ell$.  Lagrange's theorem and oddness of $q$ give
$2\ell\mid q-1$.  Two distinct factors therefore contribute at least
$(2\ell+1)^2$ to the radical.  Since $L_\ell=1$ and
$\rad(M_\ell)E_\ell=M_\ell=\Phi_\ell(2)$, the excess form follows.  If one
factor is at least $H$, combine it with the distinct factor at least
$2\ell+1$.
\end{proof}

Writing $P(N)$ for the greatest prime factor of $N$, Erd\H{o}s and Shorey prove
$P(2^\ell-1)\gg\ell\log\ell$ for prime $\ell$
\cite{ErdosShorey1976}; the later account of Ford, Luca, and Shparlinski
records the same uniform input \cite{FordLucaShparlinski2009}.  The asymmetric
part of Theorem~\ref{thm:mersenne-prime-square-carrier} therefore yields
\begin{equation}\label{eq:mersenne-literature-saving}
 \frac{E_\ell}{\Phi_\ell(2)}
       \ll\frac1{\ell^2\log\ell}
\end{equation}
for every sufficiently large prime $\ell$, including the prime
$M_\ell$ case trivially.  This is only a polynomial saving.  It leaves
\[
 \log E_\ell\le \ell\log2-2\log\ell-\log\log\ell+O(1),
\]
whereas the Mersenne endpoint needs $\log E_\ell=o(\ell)$.

The corrected global ledger identifies a clean sufficient target.

\begin{proposition}[Base-block totient criterion]
\label{prop:mersenne-totient-criterion}
If
\[
                         \log E_d=o(\varphi(d)),
\]
then $\log W_m=o(m)$ as $m\to\infty$.
\end{proposition}

\begin{proof}
Put $e_d=\log E_d\ge0$.  Given $\eps>0$, choose $D$ so that
$e_d\le\eps\varphi(d)$ for $d\ge D$, and put
$C_D=\sum_{d<D}e_d$.  The divisor identity
$\sum_{d\mid m}\varphi(d)=m$ gives
\[
 \log B_m=\sum_{d\mid m}e_d\le C_D+\eps m.
\]
Theorem~\ref{thm:mersenne-order-decomposition} now gives
$\log W_m\le\log m+C_D+\eps m$.  Divide by $m$, let $m$ tend to infinity,
and then let $\eps$ tend to zero.
\end{proof}

\begin{proposition}[Cyclotomic reformulation]
\label{prop:mersenne-cyclotomic-reformulation}
For every $d>1$,
\begin{equation}\label{eq:mersenne-cyclotomic-excess}
 E_d=\frac{\Phi_d(2)}{\rad(\Phi_d(2))},\qquad
 \log\Phi_d(2)=\varphi(d)\log2+O(1),
\end{equation}
where the absolute value of the $O(1)$ term is at most two.  If $V_d$ is
the powerful part of $\Phi_d(2)$, then
\[
                         E_d\le V_d\le E_d^2.
\]
Hence the hypothesis of
Proposition~\ref{prop:mersenne-totient-criterion} is equivalent both to
\[
 \log\rad(\Phi_d(2))=\varphi(d)\log2+o(\varphi(d))
\]
and to $\log V_d=o(\varphi(d))$.
\end{proposition}

\begin{proof}
Murty and Wong record the unconditional cyclotomic classification that a
prime divisor of $\Phi_d(a,b)$ is $1\pmod d$, except possibly for the largest
prime factor of $d$, and that the exceptional divisor occurs to at most the
first power \cite{MurtyWong2002}.  Pomerance gives the same classification
directly for $\Phi_d(2)$ and records the sharper bounds
$2^{\varphi(d)-1}\le\Phi_d(2)<2^{\varphi(d)+1}$
\cite{PomeranceCyclotomic2025}.  For $(a,b)=(2,1)$ the first class consists
of exact-order-$d$ primes; the exceptional class disappears on division by
the radical.  This proves the first identity in
\eqref{eq:mersenne-cyclotomic-excess}.  The cyclotomic product formula gives
\[
 \log\Phi_d(2)=\varphi(d)\log2+
   \sum_{e\mid d}\mu(d/e)\log(1-2^{-e}).
\]
Since $-\log(1-x)\le2x$ for $0\le x\le1/2$, the last sum has absolute value
at most $\sum_{e\ge1}2^{1-e}=2$.  A prime power $q^a$, $a\ge2$, contributes
$q^{a-1}$ to $E_d$ and $q^a$ to $V_d$; now
$a-1\le a\le2(a-1)$.  Taking logarithms proves the equivalences.
\end{proof}

The reformulation permits an unconditional removal of all sufficiently
small repeated-prime support.  Put
\[
 \mathcal W_d=\{q:\operatorname{ord}_q(2)=d,
                    \ v_q(2^d-1)\ge2\},
\]
\[
 T_d=\prod_{q\in\mathcal W_d}q,\qquad
 D_d=\prod_{\substack{\operatorname{ord}_q(2)=d\\v_q(2^d-1)\ge3}}
              q^{v_q(2^d-1)-2}.
\]
Thus $E_d=T_dD_d$.

\begin{proposition}[Near-quadratic tail reduction]
\label{prop:mersenne-near-quadratic-tail}
Set
\[
 Y_d=\frac{\varphi(d)^2}{\log\log(3d)},\qquad
 T_d^{\le Y}=\prod_{\substack{q\in\mathcal W_d\\q\le Y_d}}q.
\]
Then
\begin{equation}\label{eq:mersenne-small-wieferich-mass}
                   \log T_d^{\le Y}=o(\varphi(d)).
\end{equation}
If $\log E_d=o(\varphi(d))$ fails, then there are some $\epsilon>0$ and a
sequence $d_j\to\infty$ such that, for each fixed $\delta>0$, there is a
subsequence on which at least one of the following persists:
\begin{enumerate}
\item $\log D_{d_j}\ge\epsilon\varphi(d_j)/4$;
\item the interval $Y_{d_j}<q\le d_j^{2+\delta}$ contains at least
\[
 \frac{\epsilon\varphi(d_j)}{4(2+\delta)\log d_j}
\]
distinct base-two Wieferich primes, all of exact order $d_j$;
\item
\[
 \sum_{\substack{\operatorname{ord}_q(2)=d_j,
                  \ v_q(2^{d_j}-1)\ge2\\q>d_j^{2+\delta}}}
       \log q\ge\epsilon\varphi(d_j)/4.
\]
Every prime in the third arm has
$\operatorname{ord}_q(2)<q^{1/(2+\delta)}$ and hence belongs to the
zero-density exceptional set in the Erd\H{o}s--Murty almost-all order
theorem \cite{ErdosMurty1999}.  That density theorem does not control the
weighted mass on a sparse sequence.
\end{enumerate}
\end{proposition}

\begin{proof}
Exact order $d$ implies $q\equiv1\pmod d$.  Fix
$\theta\in(1/2,1)$.  The uniform bound
$\varphi(d)\gg d/\log\log(3d)$ gives $d<Y_d^\theta$ for all sufficiently
large $d$.  The weighted Brun--Titchmarsh inequality of Murty--S\'eguin
\cite[Theorem~2.4]{MurtySeguin2019} then gives
\[
 \log T_d^{\le Y}
 \le\sum_{\substack{q\le Y_d\\q\equiv1\pmod d}}\log q
 \le\frac{2Y_d\log Y_d}
          {\varphi(d)\log(Y_d/d)}
 =o(\varphi(d)),
\]
which proves \eqref{eq:mersenne-small-wieferich-mass}.  If the target fails,
choose $d_j$ with
$\log E_{d_j}\ge\epsilon\varphi(d_j)$.  Split
$E_d=T_dD_d$ into the small factor, the two displayed support ranges, and
$D_d$.  For large $j$ the small logarithm is at most
$\epsilon\varphi(d_j)/4$; one of the other three nonnegative logarithms is
therefore at least $\epsilon\varphi(d_j)/4$.  Infinite pigeonhole fixes an
arm.  In the transition range every summand is at most
$(2+\delta)\log d_j$, giving the asserted cardinality.  In the extreme range
$q>d_j^{2+\delta}$ and exact order give
$d_j<q^{1/(2+\delta)}$.  The final literature comparison follows from the
almost-all lower order bound of Erd\H{o}s and Murty.
\end{proof}

\begin{corollary}[A fixed fraction of the square budget]
\label{cor:mersenne-transition-square-budget}
Let $s_d=|\mathcal W_d|$ and put
\[
 B_d=\frac{(\varphi(d)+1)\log2}{2\log(d+1)}.
\]
Then $s_d\le B_d$.  On the transition alternative of
Proposition~\ref{prop:mersenne-near-quadratic-tail},
\[
 \frac{|\{q\in\mathcal W_d:Y_d<q\le d^{2+\delta}\}|}{B_d}
 \ge
 \frac{\epsilon}{2(2+\delta)\log2}
 \frac{\varphi(d)}{\varphi(d)+1}
 \frac{\log(d+1)}{\log d}.
\]
Thus that alternative consumes a fixed positive proportion of the maximum
support allowed by square divisibility alone.
\end{corollary}

\begin{proof}
Every $q\in\mathcal W_d$ is at least $d+1$, while
$T_d^2\mid\Phi_d(2)$.  Pomerance's bound
$\Phi_d(2)<2^{\varphi(d)+1}$ \cite{PomeranceCyclotomic2025} therefore gives
\[
 s_d\log(d+1)\le\log T_d
 \le\tfrac12\log\Phi_d(2)
 <\tfrac12(\varphi(d)+1)\log2.
\]
Divide the transition-cardinality bound in the preceding proposition by
$B_d$.
\end{proof}

This is a necessary structure, not a contradiction.  One very deep lift can
still carry the first arm, and no unconditional theorem excludes the
fixed-proportion cluster in the second arm.

Because $q\nmid(q-1)/d_q$, LTE also gives
\[
 v_q(2^{q-1}-1)=v_q(2^{d_q}-1)=w_q.
\]
Thus $T_d$ is exactly the same-order Wieferich support and $D_d$ is its
super-Wieferich excess.  A prime enters only one canonical block.  Finiteness
of base-two super-Wieferich primes would remove the first arm eventually but
would leave the two support tails.  Current fixed-base results, including
Fellini--Murty's finite-super implication, do not provide a weighted estimate
inside one exact-order block \cite{FelliniMurty2026}.  Under uniformly bounded
canonical valuations---in particular, under finiteness of base-two
super-Wieferich primes---Murty--S\'eguin's bound for the largest prime is only
at a polynomial scale
\cite{MurtySeguin2019}.  Pomerance's
unconditional results establish many composite primitive parts; his stronger
distinct-factor theorem is
abc-conditional \cite{PomeranceCyclotomic2025}.  Neither conclusion yields
exponential radical saturation on the $\varphi(d)$ scale.

\begin{proposition}[A nontrivial canonical block]
\label{prop:mersenne-1093-block}
The canonical exact-order block at $d=364$ is nontrivial; more precisely,
\[
                         1093\mid E_{364}.
\]
\end{proposition}

\begin{proof}
Direct integer arithmetic gives
\[
 1093^2\mid2^{364}-1,\qquad 1093^3\nmid2^{364}-1,
\]
and
\[
 2^{182}\equiv1092,\qquad 2^{52}\equiv27,\qquad
 2^{28}\equiv121\pmod {1093}.
\]
The number $1093$ is prime and $364=2^2\cdot7\cdot13$.  The prime-divisor
criterion for multiplicative order therefore gives
$\operatorname{ord}_{1093}(2)=364$: exponent $364$ kills $2$, while division
by any prime divisor of $364$ does not.  The first two divisibilities show
that the valuation at $1093$ is exactly two.  Thus the defining product for
$E_{364}$ contains $1093^{2-1}$.
\end{proof}

Proposition~\ref{prop:mersenne-1093-block} is a complete counterexample to
the strengthening $E_d=1$ for every $d$.  It does not contradict
$\log E_d=o(\varphi(d))$, because a single fixed block has no effect on the
limit at infinity.

Two further exact witnesses delimit tempting universal strengthenings:
\[
 2^{37}-1=223\cdot616318177
\]
has exactly two prime factors, refuting the claim that every composite prime
layer with $\ell\ge37$ has at least three; and
\[
 2^{11}-1=23\cdot89,\qquad
 \rad(M_{11})=2047<(2\cdot11+1)^3
\]
refutes a cubic radical replacement valid at every composite odd-prime
layer.  These witnesses do not refute an eventual three-factor estimate, a
weighted base-block theorem, or Proposition~\ref{prop:mersenne-totient-criterion}.
A deterministic scan of all prime indices $3\le\ell\le61$ found no repeated
factor; that is recorded only as a finite no-hit.

Lean verifies the elementary prime-layer theorems, the three prime-layer
witnesses, and the finite totient-tail inequality.  Companion modules additionally
verify exact-order LTE, $L_m\mid m$, the finite product
$W_m=L_m\prod_{d\mid m}E_d$, equality of relative and canonical blocks on
divisors, the exact logarithmic divisor-sum identity, and the conditional
asymptotic conclusion of Proposition~\ref{prop:mersenne-totient-criterion},
as well as the nontrivial canonical-block witness $1093\mid E_{364}$.
The finite core of the powerful-part comparison, the deep/transition/extreme
mass trichotomy, and both the realized-support and exact ambient square-budget
ratio implications is checked in a
fifth module.  That module takes the cyclotomic identification, the
Brun--Titchmarsh small-arm estimate, and the prime-order distribution theorem
as explicit paper inputs; it does not assert them as Lean theorems.
The cited largest-prime-factor theorem remains an external published input,
and $\log E_d=o(\varphi(d))$ remains an open hypothesis.  The lifting
remainder and the passage from that hypothesis are now kernel-checked.
